\subsection{Index of a Fredholm map}

Let E, F, and G be locally convex spaces.

Let \(u \colon \mathrm{E} \to \mathrm{F}\) be a Fredholm operator. The vector spaces \(\operatorname{Ker}(u)\) and \(\operatorname{Coker}(u)\) are of finite dimension (Lemma 1 of III, p. 41). Recall that the integer

\[
\dim \operatorname{Coker} (u) - \dim \operatorname{Ker} (u) = \operatorname{codim} _ {\mathrm{F}} \operatorname{Im} (u) - \dim \operatorname{Ker} (u) \tag {1}
\]

is called the index of u and is denoted ind(u) (A, V, p. 126).

If \(u \colon \mathrm{E} \to \mathrm{F}\) and \(v \colon \mathrm{F} \to \mathrm{G}\) are Fredholm maps, then so is \(v \circ u\) (III, p. 41, remark 2), and one has (A, V, p. 127, lemma 2)

\[
\operatorname{ind} (v \circ u) = \operatorname{ind} (v) + \operatorname{ind} (u).\tag{2}
\]

Suppose E and F are separated, and let \(u\colon\mathrm{E}\to\mathrm{F}\) be a Fredholm operator; let us adopt the notations of condition (iv) of Prop. 2 of III, p. 42. We then have \(\mathrm{ind}(u)=\mathrm{dim}(\mathrm{F}_{2})-\mathrm{dim}(\mathrm{E}_{2})\) . Equip the dual of each of these spaces with the weak topology (resp. the compact-open topology, the topology of bounded convergence). Then \(\mathrm{E}^{\prime}\) is identified with the topological direct sum of \(\mathrm{E}_{1}^{\prime}\) and \(\mathrm{E}_{2}^{\prime}\) , and \(\mathrm{F}^{\prime}\) is identified with that of \(\mathrm{F}_{1}^{\prime}\) and \(\mathrm{F}_{2}^{\prime}\) , and \(^{t}u\) induces an isomorphism of \(\mathrm{F}_{1}^{\prime}\) onto \(\mathrm{E}_{1}^{\prime}\) and vanishes on \(\mathrm{F}_{2}^{\prime}\) . Thus the transpose \(^{t}u\colon\mathrm{F}^{\prime}\to\mathrm{E}^{\prime}\) is a Fredholm map (loc. cit.). The kernel of \(^{t}u\) is \(\mathrm{F}_{2}^{\prime}\) , and its dimension is that of \(\mathrm{F}_{2}\) , that is, that of the cokernel of \(u\) . Hence

\[
\operatorname{ind} (u) = \dim \operatorname{Ker} (^ {t} u) - \dim \operatorname{Ker} (u).\tag{3}
\]

Moreover, the image of \(^{t}u\) is \(E_{1}^{\prime}\) , and the dimension of the cokernel of \(^{t}u\) is therefore equal to the dimension of \(E_{2}^{\prime}\) , which is that of the kernel \(E_{2}\) of u. We deduce from this

\[
\operatorname{ind} (^ {t} u) = - \operatorname{ind} (u).\tag{4}
\]

Suppose E and F are separated, and let \(u\colon\mathrm{E}\to\mathrm{F}\) be a Fredholm operator of index 0. Then \(u\) is a strict morphism by prop. 2 of III, p. 42. Since \(\dim\operatorname{Ker}(u)=\operatorname{codim}_{\mathrm{F}}\operatorname{Im}(u)\) , the map \(u\) is an isomorphism as soon as it is injective or surjective.

PROPOSITION 3. --- Let E and F be locally convex spaces and \(u \in \mathcal{L}(\mathrm{E};\mathrm{F})\) . Let \(\mathrm{E}_{1}\) (resp. \(\mathrm{F}_{1}\) ) be a closed subspace of finite codimension of E (resp. F). Suppose that u maps \(\mathrm{E}_{1}\) into \(\mathrm{F}_{1}\) and denote by \(u_{1} \in \mathcal{L}(\mathrm{E}_{1};\mathrm{F}_{1})\) the map which coincides with u on \(\mathrm{E}_{1}\) . For u to be a Fredholm operator, it is necessary and sufficient that \(u_{1}\) be one. We then have

\[
\operatorname{ind} (u) - \operatorname{ind} \left(u _ {1}\right) = \operatorname{codim} _ {\mathrm{F}} \left(\mathrm{F} _ {1}\right) - \operatorname{codim} _ {\mathrm{E}} \left(\mathrm{E} _ {1}\right).\tag{5}
\]

Let \(i\colon \mathrm{E}_{1}\to \mathrm{E}\) and \(j\colon \mathrm{F}_{1}\to \mathrm{F}\) be the canonical injections. These are Fredholm maps, and one has

\[
\operatorname{ind} (i) = \operatorname{codim} _ {\mathrm{E}} (\mathrm{E} _ {1}), \quad \operatorname{ind} (j) = \operatorname{codim} _ {\mathrm{F}} (\mathrm{F} _ {1}).
\]

Since \(j \circ u_{1} = u \circ i\) , one sees that u is a Fredholm map if and only if \(u_{1}\) is one (remarks 3 and 4 of III, p. 41). If this is so, we have

\[
\operatorname{ind} (j) + \operatorname{ind} (u _ {1}) = \operatorname{ind} (j \circ u _ {1}) = \operatorname{ind} (u \circ i) = \operatorname{ind} (u) + \operatorname{ind} (i),
\]

whence formula (5).

PROPOSITION 4. --- Let E and F be separated locally convex spaces, \(u: \mathrm{E} \to \mathrm{F}\) a Fredholm operator, and \(\widehat{u}: \widehat{\mathrm{E}} \to \widehat{\mathrm{F}}\) the extension of u to the completions. Then \(\widehat{u}\) is a Fredholm operator, and one has \(\operatorname{Ker}(\widehat{u}) = \operatorname{Ker}(u)\) and \(\operatorname{ind}(\widehat{u}) = \operatorname{ind}(u)\) .

Let us adopt the notation of condition (iv) of Prop. 2 of III, p. 42. Since the vector spaces \(E_{2}\) and \(F_{2}\) are finite-dimensional, they are complete. The completion of \(E_{1}\) (resp. \(F_{1}\) ) is identified with the closure of \(E_{1}\) in \(\widehat{E}\) (resp. of \(F_{1}\) in \(\widehat{F}\) ), and \(\widehat{E}\) (resp. \(\widehat{F}\) ) is the topological direct sum of \(\widehat{E}_{1}\) and \(E_{2}\) (resp. \(\widehat{F}_{1}\) and \(F_{2}\) ). The linear map \(\widehat{u}\) defines by restriction an isomorphism of \(\widehat{E}_{1}\) onto \(\widehat{F}_{1}\) and vanishes on \(E_{2}\) . The proposition then follows from implication (iv)⇒(i) of loc. cit.

